%% Chapter 13: Approximation Algorithms II
%% Status: skeleton -- learning goals and section outline fixed, prose to be written.

\chapter{Approximation Algorithms II}
\label{ch:approx-ii}

This chapter adds two things to the previous one: linear programming, which
turns the search for a lower bound on $\OPT$ into a mechanical procedure, and
approximation schemes, which let us trade running time against accuracy. We also
see the first limits --- a problem can be constant-factor approximable and still
provably admit no FPTAS.

\section*{Learning goals}
\addcontentsline{toc}{section}{Learning goals}

After this chapter you should be able to:
\begin{itemize}[leftmargin=*]
  \item formulate an LP relaxation of a combinatorial problem and round it;
  \item explain the primal--dual method;
  \item define PTAS and FPTAS and give an FPTAS for \lang{Knapsack};
  \item prove that the shortest-path-metric MST heuristic is a $2$-approximation for \lang{Steiner Tree};
  \item prove that a strongly NP-hard problem with integral objective admits no FPTAS unless P${}={}$NP.
\end{itemize}

\section{LP relaxation and rounding}
\label{sec:lp-rounding}

\todo{Write this section.}

\section{The primal--dual method}
\label{sec:primal-dual}

\todo{Write this section.}

\section{Approximation schemes}
\label{sec:approximation-schemes}

\todo{Write this section.}
% Planned content: PTAS, EPTAS, FPTAS, and an FPTAS for \lang{Knapsack} by rounding the profits.

\section{Steiner tree}
\label{sec:steiner-tree}

\todo{Write this section.}
% Planned content: The metric closure on the terminals, the MST heuristic, and the proof that it is a $2$-approximation.

\section{Limits: no FPTAS for Steiner tree}
\label{sec:no-fptas}

\todo{Write this section.}
% Planned content: Strong NP-hardness and gap arguments; the weaker statement about additive constants.

\section{Exercises}
\label{sec:ex-approx-ii}

\todo{Write this section.}

\begin{poolquestions}
  \poolq{14}{Approximating Steiner Tree}Part 1 in \cref{sec:steiner-tree}, part 2 in \cref{sec:no-fptas}.
\end{poolquestions}
